\documentclass[a4paper]{article}

\usepackage{graphicx}
\usepackage{amsmath,amssymb}
\usepackage{minted}
\usepackage{multirow}
\usepackage{float}
\usepackage{algorithm}
\usepackage{algpseudocode}
\usepackage{todonotes}
\usepackage{forest}
\usepackage{xstring}
\usepackage{xcolor}
\usepackage[most]{tcolorbox}
\usepackage{xparse}
\usepackage{natbib}
\usepackage{adjustbox}

\usetikzlibrary{graphs,graphdrawing}
\usegdlibrary{layered}

% for external graphviz
\input{/home/donkarlo/Dropbox/repo/nd_language_ai_project/src/nd_language_ai/formal/computer/programming/tex/latex/code_clips/external_graphviz.tex}

% for tags
\input{/home/donkarlo/Dropbox/repo/nd_language_ai_project/src/nd_language_ai/formal/computer/programming/tex/latex/part/control_sequence/kind/semtag/semtag.tex}

% for links
\input{/home/donkarlo/Dropbox/repo/nd_language_ai_project/src/nd_language_ai/formal/computer/programming/tex/latex/code_clips/seehere.tex}

% for nested lists and sections
\input{/home/donkarlo/Dropbox/repo/nd_language_ai_project/src/nd_language_ai/formal/computer/programming/tex/latex/code_clips/deep_structure.tex}

\title{}
\date{}
\author{Mohammad Rahmani}

\begin{document}
   \maketitle
   \begin{table}[H]

   \centering

   \small

   \setlength{\tabcolsep}{4pt}

   \begin{tabular}{p{0.25\textwidth} p{0.18\textwidth} p{0.47\textwidth}}

      \hline

      \textbf{Concept} & \textbf{Variable name} & \textbf{Description} \\

      \hline

      State neuron
      & $X_{i,\ell}^{(d)}$
      & State neuron identified by neuron index $i$, representation level $\ell$, and degree $d$. \\

      Neuron identifier
      & $i$
      & Identifier of an individual neuron. \\

      Representation level
      & $\ell$
      & Level at which the state neuron represents information, where $\ell \in L$. \\

      Representation levels
      & $L$
      & Set of representation levels: $L=\{\text{point},\text{phoneme},\text{word},\text{sentence},\text{final concept}\}$. \\

      Degree
      & $d$
      & Determines the represented state quantity. For $d=0$, the neuron represents intensity; for $d=2$, it represents the intensity change rate. \\

      Sensory receptor neuron
      & $S_{i}^{(d)}$
      & Sensory receptor neuron identified by neuron index $i$ and degree $d$. \\

      Observation neuron
      & $Z_{i}^{(d)}$
      & Observation neuron identified by neuron index $i$. \\

      \hline

   \end{tabular}

   \caption{Notation for state neurons, sensory receptor neurons, and observation neurons.}

   \label{tab:neuron-state-observation-notation}

\end{table}
   \bibliographystyle{plainnat}
   \bibliography{/home/donkarlo/Dropbox/repo/nd_shared_research/bibliography/bibliography.bib}
\end{document}
